% =============================================================
%  Module Description — Probability Theory and Statistics (MAT202-S28)
%  Shown on https://wosm.academy as the course page for MAT202-S28.
%  Compile with: pdflatex mat202-module-description.tex
%  Write "TODO(owner): ..." where information is still missing;
%  the website highlights it.
% =============================================================
\documentclass[11pt,a4paper]{article}

\usepackage[utf8]{inputenc}
\usepackage[T1]{fontenc}
\usepackage{lmodern}
\usepackage[margin=2.2cm]{geometry}
\usepackage{amssymb}
\usepackage{booktabs}
\usepackage{tabularx}
\usepackage{array}
\usepackage{enumitem}
\usepackage{xcolor}
\usepackage{titlesec}
\usepackage{fancyhdr}
\usepackage[hidelinks]{hyperref}

% ---------- Module details (edit these) ----------------------
% \ModuleCode must equal the course's `code` in src/content/courses/.
\newcommand{\ModuleTitle}{Probability Theory and Statistics}
\newcommand{\ModuleCode}{MAT202-S28}
\newcommand{\Programme}{Foundation programme}
\newcommand{\Level}{Semester 4}
\newcommand{\Semester}{Spring 2028}
\newcommand{\Credits}{8 ECTS}
\newcommand{\ModuleLeader}{Dr. Ali Khudiyev}
\newcommand{\Email}{ali.khudiyev@wosm.academy}
\newcommand{\OfficeHours}{Monday--Friday, 10:00--17:00 (Baku time, UTC+4)}
\newcommand{\Prerequisites}{MAT102-S27 Linear Algebra; MAT231-F27 Calculus II}

% ---------- Styling ------------------------------------------
\definecolor{accent}{RGB}{0,70,130}
\titleformat{\section}{\large\bfseries\color{accent}}{\thesection.}{0.5em}{}
\titlespacing*{\section}{0pt}{1.2em}{0.5em}
\setlist{itemsep=2pt, topsep=4pt}
\setlength{\parindent}{0pt}
\setlength{\parskip}{4pt}

\pagestyle{fancy}
\fancyhf{}
\lhead{\small \ModuleCode\ -- \ModuleTitle}
\rhead{\small Module Description}
\cfoot{\small \thepage}
\renewcommand{\headrulewidth}{0.4pt}

% =============================================================
\begin{document}

\begin{center}
  {\LARGE\bfseries\color{accent} \ModuleTitle}\\[4pt]
  {\large Module Description}
\end{center}
\vspace{6pt}

% ---------- Summary table (not shown on the website: the course page
% shows these details from the macros above) --------------------------
\begin{tabularx}{\textwidth}{>{\bfseries}l X >{\bfseries}l X}
\toprule
Module code   & \ModuleCode    & Credits      & \Credits \\
Programme     & \Programme     & Level        & \Level \\
Semester      & \Semester      & Prerequisites & \Prerequisites \\
Module leader & \ModuleLeader  & Email        & \href{mailto:\Email}{\Email} \\
Office hours  & \multicolumn{3}{l}{\OfficeHours} \\
\bottomrule
\end{tabularx}

% ---------- Sections (shown on the website) ------------------
\section{Module Overview}
Probability Theory and Statistics introduces the mathematics of uncertainty and
of learning from data. The first part builds probability theory from axioms;
the second part uses it to draw conclusions from data with estimation, confidence
intervals, hypothesis tests and regression.

The module covers:
\begin{itemize}
  \item Probability spaces, conditional probability, random variables and their
        distributions (first 5--6 weeks)
  \item Expectation, variance, joint distributions, tail bounds and the limit
        theorems
  \item Statistical inference: estimation, confidence intervals, hypothesis testing
        and linear regression
  \item Computational methods: simulation, Monte Carlo and the bootstrap
\end{itemize}

Examples come from computer science throughout: randomised algorithms and hashing,
the analysis of average-case running time, A/B testing, and the probabilistic
foundations of machine learning. Students run simulations and analyse data in
short programming assignments.

\section{Aims}
\begin{itemize}
  \item Give a rigorous introduction to probability theory, using the calculus and linear algebra of earlier modules.
  \item Develop the probabilistic reasoning needed to design and analyse randomised algorithms and systems.
  \item Introduce the main methods of statistical inference and the assumptions behind them.
  \item Prepare students for machine learning, data science, information theory and networking.
\end{itemize}

\section{Learning Outcomes}
On successful completion of this module, students will be able to:
\begin{enumerate}[label=LO\arabic*.]
  \item Model random experiments with probability spaces and compute probabilities using counting, conditional probability, independence and Bayes' theorem.
  \item Work with standard discrete and continuous distributions and compute expectations, variances and distributions of functions of random variables.
  \item Handle joint distributions, covariance and conditional expectation, including the multivariate normal distribution.
  \item Apply linearity of expectation, indicator variables and the Markov, Chebyshev and Chernoff bounds to analyse randomised algorithms.
  \item State and apply the law of large numbers and the central limit theorem.
  \item Construct estimators by the method of moments and maximum likelihood, assess their bias and variance, and build confidence intervals.
  \item Carry out and interpret hypothesis tests and simple and multiple linear regression, and recognise common misuses of $p$-values.
  \item Use simulation, Monte Carlo methods and the bootstrap to approximate probabilities and quantify uncertainty in code.
\end{enumerate}

\section{Syllabus}
\begin{tabularx}{\textwidth}{c X}
\toprule
\textbf{Unit} & \textbf{Topic} \\
\midrule
1  & Sample spaces, events, the axioms of probability \\
2  & Counting and equally likely outcomes; birthday problem \\
3  & Conditional probability; law of total probability; Bayes' theorem \\
4  & Independence; randomised algorithms: verifying matrix products, min-cut \\
5  & Discrete random variables; Bernoulli, binomial, geometric, Poisson \\
6  & Expectation; linearity of expectation; indicator variables \\
7  & Variance and moments; coupon collector; balls into bins and hashing \\
8  & Continuous random variables; uniform, exponential, normal \\
9  & Functions of random variables; transformations of densities \\
10 & Joint distributions; marginal and conditional distributions \\
11 & Covariance and correlation; the multivariate normal distribution \\
12 & Conditional expectation; law of total expectation; quicksort analysis \\
13 & Moment generating functions \\
14 & Tail bounds: Markov, Chebyshev, Chernoff \\
15 & Laws of large numbers \\
16 & The central limit theorem \\
17 & Simulation and Monte Carlo methods \\
18 & Introduction to Markov chains \\
19 & Descriptive statistics and exploratory data analysis \\
20 & Point estimation: method of moments, maximum likelihood \\
21 & Bias, variance, mean squared error, consistency \\
22 & Confidence intervals; the $t$-distribution \\
23 & Hypothesis testing: $z$- and $t$-tests, $p$-values, errors of type I and II \\
24 & Chi-squared tests; A/B testing; multiple testing \\
25 & Linear regression: least squares, inference on coefficients \\
26 & The bootstrap; introduction to Bayesian inference \\
\bottomrule
\end{tabularx}

\section{Teaching and Learning Methods}
The module is self-paced. All materials are provided on Moodle: lecture notes and
recorded lectures for each unit, tutorial problem sheets with worked solutions,
short quizzes, and the briefs for the two projects. Problem sheets are for practice
and are not marked. Questions are discussed in the module channel on Discord and in
office hours. Independent study includes reading, problem sheets and project work.

\begin{tabularx}{0.6\textwidth}{X r}
\toprule
\textbf{Activity} & \textbf{Hours} \\
\midrule
Lectures            & 52 \\
Labs / Tutorials    & 38 \\
Independent study   & 150 \\
\midrule
\textbf{Total}      & \textbf{240} \\
\bottomrule
\end{tabularx}

\section{Assessment}
Knowledge is assessed with weekly unit quizzes, two projects and a final
exam. Working with simulation and real data is central to the module, so the
projects carry slightly more weight.
\begin{itemize}
  \item \textbf{Project 1 -- Simulation study:} simulate a randomised algorithm or
        hashing scheme, compare the observed behaviour with the expectation and tail
        bounds derived in the module, and illustrate the law of large numbers and
        the central limit theorem.
  \item \textbf{Project 2 -- Data analysis:} pose a question about a public data
        set, estimate parameters, build confidence intervals, carry out a hypothesis
        test or regression, and use the bootstrap; written up as a short report.
\end{itemize}
The final exam tests whether students can build a probabilistic model for a
problem, carry out the calculations, and choose and interpret an appropriate
statistical method.

\begin{tabularx}{\textwidth}{X c c c}
\toprule
\textbf{Component} & \textbf{Weight} & \textbf{Due} & \textbf{LOs} \\
\midrule
Unit quizzes (best 10 of 12) & 15\% & Weekly & LO1--LO7 \\
Project 1: Simulation study & 20\% & Week 7 & LO4, LO5, LO8 \\
Project 2: Data analysis & 25\% & Week 12 & LO6--LO8 \\
Final exam (3 hours, written) & 40\% & End of term & LO1--LO7 \\
\midrule
\textbf{Total}      & \textbf{100\%} & & \\
\bottomrule
\end{tabularx}

\section{Reading List}
\textbf{Essential}
\begin{itemize}
  \item J.\ K.\ Blitzstein, J.\ Hwang: \textit{Introduction to Probability}, 2nd ed., CRC Press, 2019. Free online at \url{https://probabilitybook.net}.
  \item L.\ Wasserman: \textit{All of Statistics: A Concise Course in Statistical Inference}, Springer, 2004.
\end{itemize}
\textbf{Recommended}
\begin{itemize}
  \item S.\ M.\ Ross: \textit{A First Course in Probability}, 10th ed., Pearson, 2019.
  \item C.\ M.\ Grinstead, J.\ L.\ Snell: \textit{Introduction to Probability}, 2nd ed., American Mathematical Society, 1997. Free online.
  \item M.\ Mitzenmacher, E.\ Upfal: \textit{Probability and Computing}, 2nd ed., Cambridge University Press, 2017.
  \item M.\ H.\ DeGroot, M.\ J.\ Schervish: \textit{Probability and Statistics}, 4th ed., Pearson, 2012.
\end{itemize}

\section{Materials and Licence}
Course materials are published under the
\href{https://creativecommons.org/licenses/by-sa/4.0/}{Creative Commons Attribution-ShareAlike 4.0}
licence unless stated otherwise. See \href{https://wosm.academy/licence/}{wosm.academy/licence} for details.

\section{Policies}
\begin{itemize}
  \item \textbf{Attendance:} Attendance is not mandatory. The module is self-paced, so you can study the materials when it suits you.
  \item \textbf{Late submissions:} Work submitted after the deadline is penalised.
  \item \textbf{Academic integrity:} All submitted work must be your own.
        See the \href{https://wosm.academy/code-of-conduct/}{Code of Conduct}.
\end{itemize}

\end{document}
